\documentclass[10pt,a4paper]{amsart}
\usepackage[margin=19mm]{geometry}
\usepackage{amsmath,amssymb,mathtools}
\usepackage[scr=rsfs,cal=euler]{mathalfa}
\usepackage{xfrac}
\usepackage[unicode,hidelinks,bookmarks=false]{hyperref}
\newcommand{\ie}{i.e.\ }
\newcommand{\cf}{cf.\ }
\newcommand{\resp}{respectively\ }
\newcommand{\Subsection}{\subsection}
\providecommand{\nobreakdash}{\nobreak\hbox{-}}
\setlength{\emergencystretch}{2em}
\multlinegap0pt
\usepackage{enumerate}
\usepackage{amssymb}
\usepackage{float}
\newtheorem{prop}{Proposition}
\title{Proofs of four generating function conjectures for arbor polytopes}
\author{Feihu Liu,Jinlong Tang}
\date{Source: arXiv:2605.08968v1 (CC BY 4.0)}
\begin{document}
\maketitle

\noindent\emph{Source and licence.} Proofs of four generating function conjectures for arbor polytopes. Version arXiv:2605.08968v1 (CC BY 4.0), distributed by arXiv under the Creative Commons Attribution 4.0 International licence, http://creativecommons.org/licenses/by/4.0/, as recorded in the arXiv licence metadata for this exact version and shown on the abstract page https://arxiv.org/abs/2605.08968v1. Full licence notice: \texttt{licenses/arxiv-2605.08968-CC-BY-4.0.txt}.\par\smallskip

\noindent\textit{Editorial scope.} Counted unit: the article's prop FormulaK\_t\_n (source0.tex lines 320--323). The article's complete original proof of it is delivered with it (source0.tex lines 324--339, one \texttt{proof} environment of 16 lines). No mathematics is written, repaired or re-derived: the statement and the proof are the article's own lines, in the article's own notation, and no retained line is substituted or re-flowed.  No further source-local context is needed: every cross-reference of every retained line resolves inside the two delivered spans.  Nothing was omitted from any retained span, so the package carries no omission record.  No retained line contains a citation, so the wrapper carries no bibliography and no bibliography entry.  No retained line contains a figure environment, an image-inclusion command or drawing code, so no image is delivered and the package declares no asset dependency.  Editorial formatting only, in addition: an AMS \texttt{amsart} preamble that restates the article's own declarations and macros used by the retained lines, namely the environments \texttt{prop} on separate counters, together with the standard packages those lines need.  The retained environments print in this document as Proposition 1 in order of appearance, whereas the article prints the same items with the numbering of its own full document.  The complete native source of the article as distributed in the arXiv e-print for this exact version is bundled unchanged as \texttt{sources/200188/source0.tex}.
\begin{prop}\label{FormulaK_t_n}
For any $n\in \mathbb{Z}_{+}$, the polynomial $K_{t_n}(X,Y)$ is given by
$$K_{t_n}(X,Y) = (1+XY)^n + \frac{XY^2}{1-Y} (1+XY)^{n-1} - \frac{XY^2}{1-Y} (Y(1+X))^{n-1}.$$
\end{prop}

\begin{proof}
For $0\le j\le k\le n$, we consider the coefficient of $X^{j}Y^{k}$ in $K_{t_{n}}(X,Y)$.

If $j=k$ which implies that the enumerated $a\in P_{t_{n}}$ possesses property $\mathrm{nz}(a)=\mathrm{ht}(a)$, so $a$ has exactly $j$ non-zero components, all of which are equal to $1$. One can easily see that there exist exactly $\binom{n}{j}$ such $a$. 

If $j<k$, it follows that the enumerated $a$ takes the value $k-j$ at the component of the root (which means $j\geq 1$), and among the remaining $n-1$ components, it possesses exactly $j-1$ non-zero entries, all of which are equal to $1$. There exist $\binom{n-1}{j-1}$ such $a$.

This yields the explicit expression for $K_{t_{n}}(X,Y)$,
\begin{align*}
 K_{t_n}(X,Y) &=\sum_{j=0}^{n}\binom{n}{j}X^{j}Y^{j}+\sum_{1\le j<k\le n}\binom{n-1}{j-1}X^{j}Y^{k}\\
 &= \sum_{j=0}^n \binom{n}{j} (XY)^j + \sum_{j=1}^{n-1} \binom{n-1}{j-1} X^j \sum_{k=j+1}^n Y^k \\
 &=(1+XY)^{n}+ \frac{XY^2}{1-Y}\sum_{j=1}^{n-1} \binom{n-1}{j-1} (XY)^{j-1} -\frac{XY^{n+1}}{1-Y}\sum_{j=1}^{n-1} \binom{n-1}{j-1} X^{j-1}\\
 &= (1+XY)^n + \frac{XY^2}{1-Y} (1+XY)^{n-1} - \frac{XY^2}{1-Y} (Y(1+X))^{n-1}.
\end{align*}
This completes the proof.
\end{proof}

\end{document}
